\chapter{The minimal mathematical toolkit}
\label{cap:strumenti}

This chapter collects the tools used throughout the rest of the document: Brownian motion, Itô's lemma, geometric
Brownian motion, change of measure, the link between expected values and differential equations, the
characteristic function and a few notions of computer arithmetic. The goal is not maximum generality but being able
to \emph{use} these tools with confidence and to understand why they work.

\section{Brownian motion}

\begin{definizione}
A standard \textbf{Brownian motion} $W_t$ is a process with $W_0=0$, continuous paths, and independent Gaussian
increments: $W_t-W_s\sim\N(0,\,t-s)$ for $s<t$.
\end{definizione}

\paragraph{Intuition: a random walk.} Split $[0,t]$ into $n$ steps of length $\Delta=t/n$ and at each step make a
jump of $\pm\sqrt\Delta$ with probability $\tfrac12$. The sum has mean $0$ and variance $n\Delta=t$; by the central
limit theorem, as $n\to\infty$ it becomes $\N(0,t)$. Two facts to remember:
\begin{itemize}
\item the typical displacement over a time $\Delta$ is $\sqrt\Delta$, much larger than $\Delta$ when $\Delta$ is
small: the paths are continuous but ``jagged'' and nowhere differentiable;
\item the \emph{square} of an increment, $(\Delta W)^2$, has mean $\Delta$ and variance $2\Delta^2$: as
$\Delta\to0$ it is practically deterministic and equal to $\Delta$. In short one writes
$(\mathrm{d}W)^2=\mathrm{d}t$.
\end{itemize}
To simulate: $W_{t+\Delta}=W_t+\sqrt\Delta\,Z$ with $Z\sim\N(0,1)$.

\section{Stochastic differential equations and Itô's lemma}
A \textbf{stochastic differential equation} (SDE) describes how a quantity changes over a time $\mathrm{d}t$:
\[
\mathrm{d}X_t=\underbrace{\mu(t,X_t)\dd t}_{\text{drift}}+\underbrace{\sigma(t,X_t)\dd W_t}_{\text{noise}} .
\]
If $X$ has these dynamics, how does a smooth function $f(t,X_t)$ change? In ordinary calculus the first-order chain
rule would be enough. Here, instead, the \emph{second}-order Taylor expansion contains
$(\mathrm{d}X)^2=\sigma^2(\mathrm{d}W)^2+\dots=\sigma^2\dd t$, which is of the same order as $\mathrm{d}t$ and cannot
be neglected.

\begin{teorema}[Itô's lemma]
\begin{equation}\label{eq:ito}
\mathrm{d}f=\Big(f_t+\mu f_x+\tfrac12\sigma^2 f_{xx}\Big)\dd t+\sigma f_x\dd W_t .
\end{equation}
\end{teorema}
The term $\tfrac12\sigma^2f_{xx}$ is the ``Itô term'': it is the correction due to the convexity of $f$ and to the
size of the fluctuations.

\section{Geometric Brownian motion (the Black model)}
The simplest model for a forward price is
\begin{equation}\label{eq:gbm}
\mathrm{d}F_t=\sigma F_t\dd W_t ,
\end{equation}
that is, Gaussian instantaneous returns $\mathrm{d}F/F$ with variance $\sigma^2\dd t$ and \emph{no drift} (the
forward is a martingale under $\Q^T$).

\paragraph{Solution, step by step.} Apply Itô to $f(F)=\ln F$: $f_F=1/F$, $f_{FF}=-1/F^2$, $\mu=0$, noise
coefficient $\sigma F$. Then
\[
\mathrm{d}\ln F_t=\Big(0+0+\tfrac12\sigma^2F^2\cdot(-\tfrac1{F^2})\Big)\dd t+\sigma F\cdot\tfrac1F\dd W_t
=-\tfrac12\sigma^2\dd t+\sigma\dd W_t .
\]
The right-hand side no longer depends on $F$, so it integrates directly:
\begin{equation}\label{eq:gbm-sol}
\ln F_T=\ln F-\tfrac12\sigma^2T+\sigma W_T
\quad\Longrightarrow\quad
F_T=F\exp\!\big(sZ-\tfrac12s^2\big),
\end{equation}
with $s=\sigma\sqrt T$ and $Z\sim\N(0,1)$.
$F_T$ is therefore \textbf{lognormal}: its logarithm is Gaussian with mean $\ln F-s^2/2$ and variance $s^2$.

\begin{perche}[title={Why does $-\tfrac12 s^2$ appear?}]
To preserve the martingale property. For a standard Gaussian, $\E[e^{sZ}]=e^{s^2/2}$ (proved by completing the
square in the integral of the density). Hence $\E[F_T]=F\,e^{-s^2/2}\,e^{s^2/2}=F$. Without the correction, the mean
of $F_T$ would grow just because the exponential is convex (Jensen's inequality). This ``convexity'' correction comes
back in every Monte Carlo scheme: get it wrong and the simulated forward acquires a spurious drift.
\end{perche}

\paragraph{Exact simulation.} From \eqref{eq:gbm-sol}, on a grid $t_0<t_1<\dots$
\begin{equation}\label{eq:log-euler}
\ln F_{t_{i+1}}=\ln F_{t_i}-\tfrac12\sigma^2\Delta_i+\sigma\sqrt{\Delta_i}\,Z_i,\qquad Z_i\ \text{iid }\N(0,1),
\end{equation}
which is \emph{exact} for any step $\Delta_i$ (the ``log-Euler'' scheme introduces no error when $\sigma$ is
constant).

\section{Correlated Brownian motions}
In stochastic-volatility models there are two sources of noise, one for the price and one for the volatility, with
correlation $\rho$: $\mathrm{d}\langle W^F,W^v\rangle_t=\rho\dd t$. To simulate them one starts from two
\emph{independent} normals $Z_v,Z_\perp$ and sets
\begin{equation}\label{eq:cholesky}
Z_F=\rho Z_v+\sqrt{1-\rho^2}\,Z_\perp .
\end{equation}
Check: $\Var Z_F=\rho^2+(1-\rho^2)=1$ and $\Cov(Z_F,Z_v)=\rho$. This is the Cholesky decomposition of the $2\times2$
correlation matrix.

\section{Change of measure (Girsanov) in one line}
Many closed forms are obtained by ``shifting'' a Gaussian. The basic fact is the following.

\begin{proposizione}[Exponential tilting]\label{prop:tilt}
If $Z\sim\N(0,1)$, then for every function $f$:
\begin{equation}\label{eq:tilt}
\E\big[e^{sZ-s^2/2}f(Z)\big]=\E\big[f(Z+s)\big].
\end{equation}
\end{proposizione}

\begin{proof}
Completing the square in the exponent, $sz-\tfrac12s^2-\tfrac12z^2=-\tfrac12(z-s)^2$, so
\[
\int_{\R} e^{sz-s^2/2}f(z)\,\frac{e^{-z^2/2}}{\sqrt{2\pi}}\dd z
=\int_{\R} f(z)\,\frac{e^{-(z-s)^2/2}}{\sqrt{2\pi}}\dd z .
\]
With the change of variable $z=y+s$ the right-hand side becomes $\E f(Z+s)$.
\end{proof}
Interpretation: weighting the scenarios with the positive, mean-one variable $L=e^{sZ-s^2/2}$ is equivalent to
changing probability, and under the new probability $Z$ has mean $s$ instead of $0$. This is the static version of
\textbf{Girsanov's theorem}: under a change of measure, a Brownian motion acquires a drift.

\paragraph{The share measure.} A case we will use often: $L=F_T/F$, which by \eqref{eq:gbm-sol} is exactly
$e^{sZ-s^2/2}$. Weighting by $F_T/F$ means using the underlying itself as numéraire. It is used to compute terms such
as $\E[F_T\ind\{F_T>K\}]$ (Chapter~\ref{cap:black}), for the symmetries of Heston and SABR and for the reflection
principle of barriers.

\section{From the expected value to the partial differential equation}
Let $v(t,F)=\E[g(F_T)\mid F_t=F]$. By the tower property, $v(t,F_t)$ is a martingale; by Itô its $\mathrm{d}t$ part
must vanish:
\begin{equation}\label{eq:feynman-kac}
\partial_tv+\tfrac12\sigma^2F^2\,\partial_{FF}v=0,\qquad v(T,F)=g(F).
\end{equation}
This is the \textbf{Feynman--Kac} link: an expected value satisfies a partial differential equation (PDE), and vice
versa. In Chapter~\ref{cap:black} this equation becomes Black's equation and explains the link between gamma and
theta; in Chapter~\ref{cap:heston} we use it to derive the Heston characteristic function.

\section{The characteristic function}
\begin{definizione}
The \textbf{characteristic function} (CF) of a real random variable $X$ is $\varphi(u)=\E[e^{iuX}]$, $u\in\R$ (and,
when the expectation exists, also for complex $u$).
\end{definizione}

\noindent Useful properties:
\begin{itemize}
\item $\varphi(0)=1$, $|\varphi(u)|\le1$ for real $u$, $\varphi(-u)=\overline{\varphi(u)}$;
\item the CF determines the distribution: the density is recovered by the inverse Fourier transform
$f(x)=\frac1{2\pi}\int e^{-iux}\varphi(u)\dd u$;
\item Gaussian $\N(m,s^2)$: $\varphi(u)=e^{ium-u^2s^2/2}$;
\item the \textbf{cumulants} $c_n$ are the coefficients of the expansion $\ln\E[e^{wX}]=\sum_nc_nw^n/n!$: $c_1$ is
the mean, $c_2$ the variance, $c_4$ measures the tails (kurtosis).
\end{itemize}

\begin{idea}
In many models (Heston among them) the density of the price at expiry is \emph{not} known in closed form, but the
characteristic function is. Fourier methods (COS, Lewis, Chapter~\ref{cap:heston}) compute prices directly from the
CF.
\end{idea}

\section{A little computer arithmetic}
The project checks formulas to $10^{-12}$ or better, so how the computer represents numbers matters.
\begin{itemize}
\item In double precision every number has about 16 significant digits. The \textbf{machine epsilon}
$\varepsilon=2^{-52}\approx{2.2\times10^{-16}}$ is the relative distance between two adjacent representable numbers
near 1. The \emph{ulp} (unit in the last place) of a number $x$ is about $\varepsilon|x|$.
\item \textbf{Cancellation:} subtracting two nearly equal numbers cancels the common digits and leaves only the
rounding error. A fundamental example: $\Phi(-8)\approx{6\times10^{-16}}$, but computing it as $1-\Phi(8)$ gives a
result that is off by about 10\%, because $\Phi(8)$ differs from 1 by only 5 ulp; beyond $d\approx8.3$ it gives
exactly 0. Remedy: compute the small tail directly, $\Phi(-d)$.
\item \textbf{Underflow:} numbers below $\approx10^{-308}$ become 0 (or lose precision even earlier). Remedy: work with
logarithms.
\item \textbf{Overflow:} above $\approx{1.8\times10^{308}}$ one gets \code{inf}. An example from the project: the
product of 250 prices around 100 is $10^{500}$; one computes the sum of the logarithms instead.
\end{itemize}

\begin{attenzione}
An implementation can be \emph{mathematically} correct and \emph{numerically} wrong. For each model the project
keeps a list of numerical pitfalls, in the form ``symptom $\to$ remedy''; in this document you will find them in
boxes like this one.
\end{attenzione}

\begin{riassunto}
\begin{itemize}
\item Brownian motion: increments $\N(0,\Delta)$, $(\mathrm{d}W)^2=\mathrm{d}t$. Itô = second-order Taylor.
\item Black: $\ln F_T\sim\N(\ln F-s^2/2,\,s^2)$; the $-s^2/2$ guarantees $\E F_T=F$. Log-Euler simulation is exact.
\item Correlation: $Z_F=\rho Z_v+\sqrt{1-\rho^2}Z_\perp$.
\item Change of measure: $\E[e^{sZ-s^2/2}f(Z)]=\E f(Z+s)$.
\item Feynman--Kac: expected value $\Leftrightarrow$ PDE. Characteristic function: known even when the density is
not.
\item Beware of cancellation, underflow and overflow.
\end{itemize}
\end{riassunto}

\begin{domande}[unbreakable]
\begin{enumerate}
\item Derive $\mathrm{d}\ln F$ with Itô for $\mathrm{d}F=\sigma F\dd W$.
\item Show that $\E[e^{sZ}]=e^{s^2/2}$. \emph{Hint:} complete the square,
$sz-\tfrac12z^2=-\tfrac12(z-s)^2+\tfrac12s^2$.
\item Why is $1-\Phi(d)$ a bad idea for $d=8$?
\end{enumerate}
\end{domande}
